\documentclass[11pt,a4paper]{article}
\usepackage[margin=1in]{geometry}
\usepackage{amsmath,amssymb,amsthm}
\usepackage{algorithm}
\usepackage{algpseudocode}
\usepackage{booktabs}
\usepackage{hyperref}

\theoremstyle{definition}
\newtheorem{definition}{Definition}[section]
\newtheorem{theorem}{Theorem}[section]
\newtheorem{lemma}[theorem]{Lemma}
\newtheorem{remark}{Remark}[section]

\title{\textbf{Rigorous Mathematical Specification, Kernel Associativity, and Linear Complexity of Linear Attention}}
\author{Team Scaibu \\ \small Email: \texttt{jaydeep.9664920749@gmail.com} \\ \small Architecture and Core Engine Team}
\date{October 2026}

\begin{document}
\maketitle

\begin{abstract}
Standard self-attention computes $\operatorname{softmax}(\mathbf{Q}\mathbf{K}^\top)\mathbf{V}$ in quadratic $\mathcal{O}(N^2)$ time due to non-linear row-wise softmax normalization. Linear Attention replaces the softmax kernel with non-negative feature maps $\phi(\mathbf{x}) = \operatorname{elu}(\mathbf{x}) + 1$. By leveraging the associativity of matrix multiplication $\phi(\mathbf{Q})(\phi(\mathbf{K})^\top \mathbf{V})$, attention evaluates in strictly linear $\mathcal{O}(N)$ time and constant memory. This specification derives kernel associativity, proves non-negativity guarantees, and details numerical verification.
\end{abstract}

\section{Mathematical Formulation}

Let feature map $\phi(x) = \operatorname{elu}(x) + 1.0 \ge 0$. The linear attention output for row $i$ is:
\begin{equation}
\mathbf{o}_i = \frac{\phi(\mathbf{q}_i)^\top \sum_{j=1}^N \phi(\mathbf{k}_j) \mathbf{v}_j^\top}{\phi(\mathbf{q}_i)^\top \sum_{j=1}^N \phi(\mathbf{k}_j) + \epsilon}
\end{equation}

\section{Complexity Analysis}
\begin{itemize}
    \item \textbf{Time Complexity:} $\Theta(N \cdot d_k \cdot d_v)$ operations instead of $\mathcal{O}(N^2 d)$.
    \item \textbf{Space Complexity:} $\Theta(d_k \cdot d_v)$ accumulator storage.
\end{itemize}

\begin{thebibliography}{9}
\bibitem{katharopoulos2020} A.~Katharopoulos et al., ``Transformers are RNNs: Fast autoregressive transformers with linear attention,'' in \emph{ICML}, 2020.
\end{thebibliography}

\end{document}
